\documentclass[12pt,a4paper,UTF8]{ctexart}

%==========================================================================
% 宏包
%==========================================================================
\usepackage{geometry}
\usepackage{amsmath,amssymb}
\usepackage{graphicx}
\usepackage{booktabs}
\usepackage{hyperref}
\usepackage{float}
\usepackage{enumitem}
\usepackage{multirow,array}
\usepackage{cite}
\usepackage{caption}
\usepackage{xcolor}

\geometry{left=3cm,right=2.5cm,top=2.5cm,bottom=2.5cm}
\hypersetup{colorlinks=true,linkcolor=blue,citecolor=blue,urlcolor=blue}

\title{\textbf{基于动态图注意力网络与多智能体强化学习的\\低空交叉路口飞行器协同管控}}
\author{彭义森\quad 指导老师：韩云祥\\[4pt]\small 四川大学}
\date{2026年7月}

\begin{document}
\maketitle

%==========================================================================
% 摘要
%==========================================================================
\begin{abstract}
城市低空交叉路口中电动垂直起降飞行器（eVTOL）的高密度安全协同管控是先进空中交通（AAM）面临的关键技术挑战。本文针对该场景下的冲突避免与流量协同优化问题，提出一种融合动态图注意力网络（DyGAT）的耦合多智能体强化学习（MARL）模型。模型包含三项核心设计：（1）DyGAT融合空间注意力、时序注意力与基于物理量（距离、速度差、航向差）的连续动态边权重，对飞行器间时空耦合关系进行细粒度建模；（2）分层奖励函数与速度自适应安全距离；（3）课程学习与基于冲突事件的episode级损失加权机制（含冲突episode的损失权重$\lambda=2.0$，无冲突episode权重$\lambda=1.0$），在不引入off-policy数据的前提下提升稀疏冲突样本的训练效率。

在21架eVTOL的二维交叉路口仿真场景（$1000\text{m}\times1000\text{m}$，含风扰与传感器噪声）中，与8种基线方法（涵盖规则调度、单智能体RL、CTDE-MARL及工程安全方法ORCA）进行了系统对比。相比性能最强的CTDE-MARL基线MAPPO，本文方法的无冲突完成率（CFCR）从$57.8\%$提升至$68.4\%$（$p<0.001$），最小间隔违反率（MSVR）从$3.0\%$降至$2.4\%$（$p=0.029$），总冲突次数（TC）从$20.3\pm2.4$降至$18.2\pm2.1$（降幅10.3\%，$p=0.180$）。消融实验表明DyGAT是性能贡献最大的单一模块，移除后TC增加173.6\%。

\noindent\textbf{关键词：}低空交通管控；多智能体强化学习；动态图注意力网络；冲突避免；集中训练分布式执行；无冲突完成率
\end{abstract}

%==========================================================================
% 1. 引言
%==========================================================================
\section{引言}

\subsection{研究背景}

城市地面交通拥堵的加剧推动了以eVTOL为核心的城市空中交通（UAM）发展。美国联邦航空管理局（FAA）预测，到2040年将有数百万架无人机与eVTOL在美国低空空域运行\cite{ref1}。城市低空交通具有高密度、强动态、高异构性三大特征——飞行器在空间中快速移动，邻域拓扑以秒级频率变化。在此背景下，交叉路口与汇合点的协同管控与冲突避免已成为低空无人机交通管理（UTM）的核心挑战。

传统空中交通流量管理（ATFM）多采用集中式优化架构，包括整数规划、蒙特卡洛仿真与基于欧拉模型的线性规划等方法\cite{ref2}。这些方法在可预测场景中效果良好，但计算复杂度随飞行器数量指数增长，难以适配低空环境的实时动态需求。模型预测控制（MPC）\cite{ref3}、最优互惠碰撞避免（ORCA）\cite{ref4}与控制屏障函数（CBF）\cite{ref5}等工程方法在形式化安全保证方面具有独特优势，但其性能高度依赖精确的环境动力学模型，且在高密度交互场景中倾向于保守。

近年来，MARL因其分布式决策与自适应学习特性被成功引入交通管控领域。Agogino与Tumer\cite{ref6,ref7}建立了基于航路点智能体的分布式ATFM学习框架；Brittain与Wei\cite{ref8}在BlueSky仿真器中验证了深度MARL的高性能冲突规避；Ghosh等\cite{ref9}提出深度集成MARL方法；Li等\cite{ref10}提出MS-GRL采用多尺度图增强RL解决密集无人机网络冲突。然而，现有工作在三个维度上仍存在不足：

\textbf{（1）时空耦合建模不足。}多数方法使用全连接网络或静态图结构编码智能体间关联。Brittain与Wei\cite{ref8}仅考虑最近邻静态关联；Spatharis等\cite{ref11}通过预定义关系矩阵编码交互；MS-GRL\cite{ref10}的多尺度图注意力虽较静态图有所改进，但其图结构仍基于固定距离阈值。这些方法的共同局限在于：当飞行器快速移动导致邻域拓扑剧烈变化时，静态或半静态图结构无法及时反映飞行器间动态变化的冲突风险等级——两架相向飞行的飞行器在距离相同时的风险远高于同向飞行，但固定阈值图无法区分这两种情形。

\textbf{（2）稀疏冲突样本利用不足。}低空交通中冲突属低概率高风险事件，训练中冲突样本占比通常低于$0.1\%$\cite{ref12,ref13}。标准mini-batch训练中冲突样本被淹没在海量正常样本中，且冲突规避的奖励信号仅在冲突发生或成功通过时才出现，导致信用分配困难。

\textbf{（3）环境不确定性建模不足。}多数研究将飞行器视为质点，忽略风扰、传感器噪声和物理运动约束\cite{ref11}，且采用固定安全距离，未考虑不同飞行速度下的差异化制动需求。

针对上述不足，本文设计了一种耦合MARL管控模型：通过DyGAT的物理量驱动动态边权重刻画时空耦合关系；通过课程学习与episode级冲突损失加权高效利用稀疏冲突样本；通过动态安全距离与物理运动约束提升环境建模的合理性。

\subsection{主要贡献}

\begin{enumerate}[label=(\arabic*)]
    \item 构建了面向低空交叉路口冲突规避的仿真环境，引入与速度正相关的动态安全距离模型、eVTOL物理运动约束（最大加速度/减速度/转弯角）及风扰与传感器噪声；
    \item 提出DyGAT网络，通过空间注意力、时序注意力与基于距离、速度差、航向差三个物理量的连续动态边权重，刻画飞行器间细粒度时空耦合关系——消融实验表明该模块的移除使冲突次数增加173.6\%；
    \item 设计课程学习与episode级冲突损失加权训练框架（$\lambda=2.0$），在保持PPO on-policy理论框架的前提下提升稀疏冲突样本的训练效率——消融实验表明该模块的移除使冲突次数增加35.2\%；
    \item 构建双层评价指标体系，区分表层指标（CR、TC、AD）与深层指标（CFCR、MCAE、MSVR），揭示通行完成率（99.1\%）与无冲突完成率（68.4\%）之间约30个百分点的显著差距；
    \item 与8种基线方法（涵盖规则调度RR/FCFS、单智能体DQN、独立MARL、CTDE-MARL四类方法学层次，以及工程安全方法ORCA）进行系统对比。
\end{enumerate}

%==========================================================================
% 2. 相关工作
%==========================================================================
\section{相关工作}

\subsection{传统ATFM与工程安全方法}

传统ATFM方法可分为集中式优化与分布式启发式两类。集中式优化包括整数规划与基于欧拉模型的线性规划\cite{ref2}，可提供最优性保证但面临可扩展性挑战。启发式方法中，Agogino与Tumer\cite{ref14}设计了基于规则的反馈控制机制，计算效率高（$<1$ms/步）但泛化至未见场景的能力有限。

工程安全方法方面，MPC\cite{ref3}通过滚动时域优化实现预测性控制；ORCA\cite{ref4}通过速度空间的几何约束实现分布式碰撞避免，计算效率高且具有形式化安全保证；CBF\cite{ref5}通过Lyapunov类约束将安全需求编码为控制输入限制。需要指出：受限于当前仿真环境的二维设定，MPC在本文的21架飞行器场景中存在计算实时性挑战（其计算时间随飞行器数量呈超线性增长），因此未纳入主实验对比。ORCA已作为工程基线纳入实验。

\subsection{强化学习在交通管控中的应用}

单智能体RL方面，DQN已被应用于单飞行器冲突规避与交叉口信号控制。Aslani等\cite{ref12}在城市地面交通信号控制中系统对比了DQN与A2C的鲁棒性。国内研究者在交通信号控制的RL应用方面也开展了大量工作\cite{ref15,ref16,ref17,ref18,ref19,ref20,ref21,ref22}。

多智能体RL方面，Ghavamzadeh等\cite{ref23}提出层级式MARL框架；Tumer与Agogino\cite{ref7,ref14,ref24}系统验证了分布式架构在ATM中的有效性；Ghosh等\cite{ref9}提出深度集成MARL方法；翟子洋等\cite{ref13}系统梳理了从独立Q-learning到QMIX、MAPPO等CTDE方法的演进路线，指出CTDE架构在多路口协同中的潜力。

Li等\cite{ref10}提出的MS-GRL与本文目标最为接近——同样面向密集低空网络的冲突解决，同样采用图注意力机制编码飞行器间交互。两者在方法学上存在三项基本差异，将在第5.1节详细讨论。

\subsection{图神经网络在交通管控中的应用}

图神经网络（GNN）因对关系型数据的自然建模能力在交通管控领域得到日益广泛的应用。MS-GRL\cite{ref10}采用多尺度图注意力网络（GAT），在100m、300m和500m三个距离尺度上分别构建子图，验证了图结构对大规模无人机冲突规避的积极作用。然而，固定距离阈值的图结构忽略了飞行器间相对运动状态的差异——距离相同但相对速度不同的飞行器对的冲突风险可能差别数个数量级。本文的DyGAT通过物理量驱动的连续动态边权重直接回应了这一局限。

\subsection{地面交通信号控制的RL研究}

地面交通信号控制领域的RL研究为本文提供了方法论参考。Aslani等\cite{ref12}验证了actor-critic架构在连续动作空间中的优势；赖建辉\cite{ref17}探索了优先经验回放在交通控制中的应用；翟子洋等\cite{ref13}对大规模交通信号控制中的RL方法进行了系统综述。尽管低空交通与地面交通在物理约束和冲突形态上存在本质差异，上述工作在训练机制和评估方法论上的经验具有可迁移性。

%==========================================================================
% 3. 问题建模与管控模型设计
%==========================================================================
\section{问题建模与管控模型设计}

本文聚焦城市低空十字交叉路口场景：$N$架eVTOL从均匀分布于圆周的不同方向飞向交叉区域中心，需在保持安全间隔的前提下以最小延迟通过交叉区域。

\subsection{分布式马尔可夫决策过程形式化}

将$N$架eVTOL的协同管控形式化定义为Dec-MDP六元组$\langle \mathcal{N}, \mathcal{S}, \mathcal{A}, \mathcal{P}, \mathcal{R}, \gamma \rangle$：

\begin{itemize}
    \item \textbf{智能体集合} $\mathcal{N} = \{1, \ldots, N\}$：训练阶段$N$由课程学习动态调整（$5\to21$），测试阶段$N=21$。
    \item \textbf{全局状态} $\mathcal{S}$：$S = (s_1, \ldots, s_N)$，$s_i = [x_i, y_i, v_i, v_i^{\text{target}}, d_i^{\text{center}}]$为飞行器$i$的五维形式化状态。需要注意：该五维向量是Dec-MDP的形式化定义变量；智能体策略网络的实际输入是经邻域扩展与历史序列编码后的局部观测$o_i$（见第3.2.1节）。全局状态$S$仅在CTDE训练阶段被GAT-Mixer用于计算全局价值函数。
    \item \textbf{联合动作} $\mathcal{A}$：$A = (a_1, \ldots, a_N)$，$a_i \in \mathbb{R}$为连续加速度指令（m/s$^2$）。
    \item \textbf{状态转移} $\mathcal{P}$：受飞行器运动学模型、风扰与传感器噪声共同影响。
    \item \textbf{联合奖励} $\mathcal{R}$：$R = \frac{1}{N}\sum_i r_i$，$r_i$由四个分量构成（第3.3节）。
    \item \textbf{折扣因子} $\gamma = 0.99$。
\end{itemize}

\subsection{环境模型}

仿真环境以二维平面模拟交叉路口空域（$1000\text{m} \times 1000\text{m}$），$N$架飞行器从圆周均匀分布的初始位置出发，航向指向中心区域。

\subsubsection{观测空间的三层定义}

为避免概念混淆，将状态/观测信息严格区分为三个层次：

\textbf{第一层——形式化状态} $s_i \in \mathbb{R}^5$：$[x_i, y_i, v_i, v_i^{\text{target}}, d_i^{\text{center}}]$，是Dec-MDP的形式化定义变量。

\textbf{第二层——局部观测} $o_i$：在$s_i$基础上进行空间扩展（邻域飞行器的相对位置$\Delta x_{ij},\Delta y_{ij}$、相对速度$\Delta v_{ij}$、加速度$a_j$、航向角$\theta_j$）和时间扩展（最近$T=10$步历史状态序列）。$o_i$是智能体网络的直接输入，维度随邻域飞行器数量$K = |\mathcal{N}_i|$动态变化。

\textbf{第三层——编码特征} $h_i^{\text{coupled}} \in \mathbb{R}^{256}$：$o_i$经LSTM编码器提取时序特征（输出64维），再经DyGAT层进行邻域信息聚合与时空耦合建模，最终得到256维耦合特征。$h_i^{\text{coupled}}$是策略头与Q值头的直接输入。

数据流为：$s_i \to o_i \to \text{LSTM} \to h_i^{\text{traj}} \to \text{DyGAT} \to h_i^{\text{coupled}} \to$ 策略/Q值输出。

\subsubsection{动作空间与物理运动约束}

动作空间为连续加速度指令$a_i \in [-a_{\max}^{\text{dec}}, a_{\max}^{\text{acc}}]$。需要诚实指出：当前主流eVTOL制造商（Joby、Volocopter、Lilium等）均未公开发布最大水平加速度的精确技术规格。学术文献中对eVTOL的典型建模假设约为2.0 m/s$^2$\cite{ref25}。本文在此基础上设定最大加速度为$3.0$ m/s$^2$，最大减速度为$4.0$ m/s$^2$——减速度高于加速度符合eVTOL安全设计惯例（紧急制动能力应高于正常加速能力），且取值的绝对值范围（2.0--4.0 m/s$^2$）覆盖了从常规运行到紧急机动的合理区间。速度范围设为$[60, 140]$ km/h，参考了当前主流eVTOL的技术水平（Volocopter约90--110 km/h，EHang 216约100--130 km/h，Joby S4巡航约240 km/h）。最大转弯角设为$\pi/6$（30$^\circ$）。速度更新模型为：

\begin{equation}
    v_i^{t+1} = \text{clip}\left(v_i^{t} + a_i \cdot \Delta t \cdot 3.6,\; 60,\; 140\right)
\end{equation}

其中$\Delta t = 0.5$ s，系数3.6实现m/s$^2$到km/h的单位转换。

\subsubsection{不确定性建模与动态安全距离}

\textbf{传感器噪声：}对位置观测添加独立同分布$\mathcal{N}(0, 0.5^2)$高斯噪声（m），仅作用于观测通道。

\textbf{风扰建模：}为每架飞行器逐时间步添加随机方向$\theta_w \sim \mathcal{U}(0, 2\pi)$、固定强度$w = 0.1$ m/s的风扰速度矢量。该模型是对真实低空风场的简化——真实风场通常使用Dryden（MIL-F-8785C）\cite{ref26}或Kaimal（IEC 61400-1）谱模型刻画，具有显著的空间相关性与时间连续性。

\textbf{动态安全距离：}冲突判定采用与空域平均速度$\bar{v}$（km/h）正相关的动态阈值：

\begin{equation}
    d_{\text{safety}} = d_{\text{base}} + \bar{v} \times 0.5
\end{equation}

其中$d_{\text{base}} = 50$ m。当$\exists i \neq j,\ \|(x_i, y_i) - (x_j, y_j)\|_2 < d_{\text{safety}}$时判定为一次冲突事件。

\subsection{分层奖励函数}

单飞行器即时奖励由四个分量构成：

\begin{equation}
    r_i = r_{\text{safe}} + r_{\text{speed}} + r_{\text{coop}} + r_{\text{eff}}
\end{equation}

\textbf{安全奖励} $r_{\text{safe}}$：三段式分段函数：
\begin{equation}
    r_{\text{safe}} =
    \begin{cases}
        -50 \exp\!\left(2\frac{d_{\text{safety}} - d_{\min}}{d_{\text{safety}}}\right), & d_{\min} < d_{\text{safety}} \\[8pt]
        -10\frac{1.2d_{\text{safety}} - d_{\min}}{0.2d_{\text{safety}}}, & d_{\text{safety}} \leq d_{\min} < 1.2d_{\text{safety}} \\[8pt]
        5, & d_{\min} \geq 1.2d_{\text{safety}}
    \end{cases}
\end{equation}
其中$d_{\min} = \min_{j \neq i} \|(x_i, y_i) - (x_j, y_j)\|_2$。

\textbf{速度奖励} $r_{\text{speed}}$：理想运行速度区间$[65, 90]$ km/h给予正奖励，与目标速度$70\pm10$ km/h和速度范围$[60, 140]$ km/h保持一致：
\begin{equation}
    r_{\text{speed}} =
    \begin{cases}
        -50 \cdot \dfrac{v_i - 90}{50}, & v_i > 90 \\[10pt]
        -30 \cdot \dfrac{65 - v_i}{5}, & v_i < 65 \\[10pt]
        20, & 65 \leq v_i \leq 90
    \end{cases}
\end{equation}

\textbf{协同奖励} $r_{\text{coop}}$：$r_{\text{coop}} = -0.2 \cdot |v_i - \bar{v}_{\text{neighbor}}|$，鼓励邻域内速度一致性。

\textbf{效率奖励} $r_{\text{eff}}$：$r_{\text{eff}} = -0.5 \cdot |v_i - v_i^{\text{target}}|$。

全局奖励取均值$R = \frac{1}{N}\sum_i r_i$。当所有$N$架飞行器均进入目标区域（距原点$<10$ m）时额外给予$+200$的完成奖励。

\subsection{动态图注意力网络（DyGAT）}

DyGAT是本文的核心技术创新，融合三种注意力机制：空间注意力（同时间步邻域关系）、时序注意力（历史轨迹演化趋势）与基于物理量的动态边权重，三者协同刻画飞行器间细粒度的时空耦合关系。

\subsubsection{空间注意力}

以LSTM编码特征$h_i^{\text{traj}} \in \mathbb{R}^{64}$为节点特征（$F_{\text{in}}=64$），在邻域半径500 m内计算注意力：

\begin{equation}
    e_{ij} = \text{LeakyReLU}_{0.2}\!\left(a^{\top}[W h_i^{\text{traj}} \parallel W h_j^{\text{traj}}]\right)
\end{equation}
\begin{equation}
    \alpha_{ij} = \frac{\exp(e_{ij})}{\sum_{k \in \mathcal{N}_i} \exp(e_{ik})},\quad
    h_i^{\text{spatial}} = \sum_{j \in \mathcal{N}_i} \alpha_{ij} \cdot W h_j^{\text{traj}}
\end{equation}

其中$a \in \mathbb{R}^{2F_{\text{out}}}$为可学习注意力向量，$W \in \mathbb{R}^{F_{\text{out}} \times 64}$，$F_{\text{out}}=256$。

\subsubsection{时序注意力}

对最近$T=10$步的历史特征序列计算时序注意力权重，输出$h_i^{\text{temporal}}$。空间与时序特征融合：

\begin{equation}
    h'_i = h_i^{\text{spatial}} + \alpha_{\text{temp}} \cdot h_i^{\text{temporal}}
\end{equation}

其中时序融合系数$\alpha_{\text{temp}} = 0.3$。实验测试了$\{0.1, 0.2, 0.3, 0.5, 0.7\}$，0.3在验证场景的冲突次数上表现最优。

\subsubsection{基于物理量的动态边权重}

飞行器间的物理交互强度具有明确的先验知识——距离越近、速度越接近、航向越一致的飞行器对，其潜在协同或冲突关系越强。DyGAT通过基于物理量直接计算的动态边权重编码这一先验：

\begin{equation}
    w_{ij} = \exp\!\left(-\frac{d_{ij}}{1000}\right) \cdot \exp\!\left(-\frac{|v_i - v_j|}{20}\right) \cdot \exp\!\left(-\frac{|\theta_i - \theta_j|}{\pi/4}\right)
\end{equation}

三个指数核的函数形式和衰减常数（1000 m、20 km/h、$\pi/4$ rad）为手工设计，权重值在$(0, 1]$之间连续变化。$w_{ij}$通过Hadamard积引导空间注意力：

\begin{equation}
    e_{ij}^{\text{final}} = e_{ij} \odot w_{ij}
\end{equation}

这一"物理先验$\times$学习注意力"的双层设计的优势在于：（1）物理先验计算无需可学习参数，不增加梯度回传开销；（2）$w_{ij} \in (0, 1]$的值域确保物理先验不会完全阻断任一邻域飞行器的信息流；（3）学习注意力系数$e_{ij}$可在物理先验基础上自适应调整。

\subsection{单智能体网络与混合网络}

\textbf{单智能体网络}包含四个组件：（1）2层LSTM编码器（64维，Dropout 0.1）提取轨迹时序特征；（2）DyGAT层进行邻域聚合，输出256维耦合特征$h_i^{\text{coupled}}$；（3）策略头——2层MLP（256$\to$256$\to$1，Tanh），输出归一化至$[-1,1]$的加速度均值$\mu_i$，实际加速度从$\mathcal{N}(\mu_i, 0.1^2)$采样；（4）Q值头——3层MLP（256$\to$256$\to$128$\to$1，ReLU）。

\textbf{GAT-Mixer混合网络}用于CTDE训练阶段的全局价值估计。将各智能体的4维局部状态$(x_i, y_i, v_i, v_i^{\text{target}})$作为图节点，在$N$个节点间构建完全图，通过GAT进行信息聚合。取各节点经GAT聚合后特征的均值得到全局表征$h_{\text{global}} \in \mathbb{R}^{256}$，再经MLP（256$\to$256$\to$1）输出$Q_{\text{tot}}$。需要说明：本文的GAT-Mixer是基于GAT注意力聚合与均值池化的混合网络，其设计目标是替代QMIX的超网络结构以实现更灵活的多智能体价值分解——它在技术路线上独立于ICLR 2023发表的GraphMixer\cite{ref28}（后者基于MLP-Mixer架构用于时序链路预测）。

\subsection{训练机制}

\subsubsection{课程学习}

训练难度从$N=5$架、$d_{\text{base}}=80$ m开始。进阶条件为连续3个评估周期同时满足冲突率低于阈值且完成率高于80\%：$N \leftarrow \min(N+2, 21)$，$d_{\text{base}} \leftarrow \max(d_{\text{base}}-5, 50)$。设有回退机制。最终覆盖$N \in \{5, 7, 9, 11, 13, 15, 17, 19, 21\}$共九个难度级别。

\subsubsection{基于冲突事件的Episode级损失加权}

本文严格保持在PPO的on-policy理论框架内——不使用经验回放缓冲区，不涉及跨策略版本的数据复用。PPO的标准裁剪代理目标为：

\begin{equation}
    L^{\text{CLIP}}(\theta) = \hat{\mathbb{E}}_t\left[\min\left(r_t(\theta)\hat{A}_t,\ \text{clip}(r_t(\theta), 1-\epsilon, 1+\epsilon)\hat{A}_t\right)\right]
\end{equation}

其中$r_t(\theta) = \pi_\theta(a_t|o_t) / \pi_{\theta_{\text{old}}}(a_t|o_t)$，$\epsilon = 0.2$，$\hat{A}_t$为GAE（$\lambda_{\text{GAE}}=0.95$）优势估计。

在此基础上，对包含冲突事件的episode施加更高的损失权重，以提升稀疏冲突样本的训练效率：

\begin{equation}
    L_{\text{total}} = \lambda \cdot (L_{\text{policy}} + L_q + 0.5 \cdot L_{\text{value}})
\end{equation}
\begin{equation}
    \lambda = \begin{cases}
        2.0, & \text{若episode内存在至少一次冲突事件} \\
        1.0, & \text{若episode全程无冲突}
    \end{cases}
\end{equation}

该机制与优先经验回放（PER）有本质区别：PER通过改变采样分布使高TD-error的transition被更频繁地回放，涉及不同策略版本间的数据混合，需要重要性采样修正；本文的$\lambda$加权仅缩放损失函数幅度（$\nabla_\theta(\lambda L) = \lambda \nabla_\theta L$），所有训练数据均来自当前策略$\pi_{\theta_{\text{old}}}$的最近一次采样，不涉及off-policy数据。需要指出：$\lambda$加权与PPO clip机制的交互效应值得进一步分析——当某transition的策略比率$r_t(\theta)$已被clip函数截断（梯度为零）时，$\lambda$对该transition的额外加权是无效的；这意味着$\lambda$加权的实际效果可能在不同训练阶段有所差异。

$\lambda=2.0$的取值通过在$\{1.0, 1.5, 2.0, 2.5, 3.0\}$五个离散值上的实验确定（见表\ref{tab:lambda}）。五个离散取值中的最优未必是连续空间中的全局最优。

\begin{table}[H]
    \centering
    \caption{$\lambda$取值实验结果（21架eVTOL，5次独立测试，均值$\pm$标准差）}
    \label{tab:lambda}
    \begin{tabular}{|c|c|c|c|}
        \hline
        $\lambda$ & TC$\downarrow$ & CFCR(\%)$\uparrow$ & 收敛轮次 \\
        \hline
        1.0 & $24.6 \pm 2.8$ & $53.2 \pm 2.3$ & $\sim800$ \\
        1.5 & $21.8 \pm 2.5$ & $58.7 \pm 2.0$ & $\sim600$ \\
        \textbf{2.0} & $\mathbf{18.2 \pm 2.1}$ & $\mathbf{68.4 \pm 1.8}$ & $\sim400$ \\
        2.5 & $19.5 \pm 2.3$ & $64.1 \pm 1.9$ & $\sim350$ \\
        3.0 & $22.3 \pm 3.1$ & $57.3 \pm 2.5$ & $\sim300$ \\
        \hline
    \end{tabular}
\end{table}

$\lambda > 2.0$时收敛更快但最终性能下降，推测原因为过度加权导致策略对早期冲突episode过拟合，限制了在高密度场景下的泛化能力。PTR-PPO\cite{ref27}等工作也探索了轨迹级优先级在on-policy算法中的应用——与本文的$\lambda$加权不同，PTR-PPO采用了更复杂的截断重要性采样机制来处理来自旧策略的off-policy轨迹。

%==========================================================================
% 4. 实验
%==========================================================================
\section{实验与结果分析}

\subsection{实验设置}

实验基于PyTorch实现，硬件环境为NVIDIA GeForce RTX 4060 GPU（8 GB显存）与Intel Core i9-14900K CPU，软件环境为Python 3.11。仿真参数汇总于表\ref{tab:env}。

\begin{table}[H]
    \centering\small
    \caption{仿真环境参数}
    \label{tab:env}
    \begin{tabular}{|c|c|c|c|}
        \hline
        \textbf{参数} & \textbf{取值} & \textbf{参数} & \textbf{取值} \\
        \hline
        空域大小 & $1000\text{m}\times1000\text{m}$ & 最大转弯角 & $\pi/6$ rad \\
        $d_{\text{base}}$ & $50$ m（初始80 m，递减至50 m） & 时间窗口$T$ & $10$步（5 s） \\
        $\Delta t$ & $0.5$ s & 邻域半径 & $500$ m \\
        最大步长 & $200$步（100 s） & 训练轮次 & $1000$ \\
        目标速度 & $70 \pm 10$ km/h & 折扣因子$\gamma$ & $0.99$ \\
        速度范围 & $[60, 140]$ km/h & PPO裁剪参数$\epsilon$ & $0.2$ \\
        最大加速度 & $3.0$ m/s$^2$ & 初始学习率$\eta$ & $5 \times 10^{-4}$ \\
        最大减速度 & $4.0$ m/s$^2$ & 学习率衰减率 & $0.99$/轮 \\
        风扰强度 & $0.1$ m/s & 传感器噪声$\sigma$ & $0.5$ m \\
        \hline
    \end{tabular}
\end{table}

\subsubsection{基线方法}

选取8种基线方法，涵盖四个类别：

\textbf{（1）传统规则调度（2种）：}RR（固定轮询调度）和FCFS（先到先服务）。

\textbf{（2）单/独立智能体RL（2种）：}集中式DQN（全局状态拼接为单一状态向量）和独立A2C（每架飞行器独立运行A2C，无显式协同机制）。

\textbf{（3）CTDE-MARL（4种）：}MADDPG、COMA、QMIX和MAPPO——后者是当前性能最强的CTDE-MARL基线之一。

\textbf{（4）工程安全方法（1种）：}ORCA\cite{ref4}\footnote{VO（速度障碍法）作为ORCA的前身，在高密度场景（$N=21$）中性能被ORCA一致主导（TC更高、CFCR更低），因此主实验仅报告ORCA的结果。}。ORCA实现采用RVO2库的Python接口（pyrvo），关键参数为：时间窗口$\tau=5.0$ s，安全半径等于环境$d_{\text{safety}}$（动态值），飞行器首选速度设定为指向交叉中心方向、大小等于目标速度（$70\pm10$ km/h）。MPC方法在本文场景中存在计算实时性挑战，未纳入主实验对比。

\subsubsection{评价指标}

\textbf{表层指标：}累计奖励（CR）、总冲突次数（TC）、平均延迟（AD）。

\textbf{深层指标：}通行完成率（PCR）——成功到达目标区域的飞行器占比；无冲突完成率（CFCR）——全程零冲突且所有飞行器完成通行的episode占比（CFCR是PCR的严格子集）；每架次平均冲突次数（MCAE）；最小间隔违反率（MSVR）——存在安全距离违反的时间步占比。

PCR与CFCR的区分是本文核心方法论贡献之一：以本文方法为例，PCR $=99.1\%$与CFCR $=68.4\%$之间约30个百分点的差值意味着约31\%的episode中，所有飞行器虽最终到达目标区域，但途中至少经历了一次安全距离违反。若仅报告PCR，读者将获得"系统近乎完美"的错误印象。

\subsection{训练收敛性分析}

模型在1000轮训练中的收敛过程可划分为三个特征阶段（见图\ref{fig:training}）：

\textbf{阶段I——快速收敛期（0--200轮）：}累计奖励从约$-2000$持续单调上升至约$+800$，冲突次数从约80次/episode快速下降至约25次/episode，平均延迟从约$1.8$ s降至约$0.8$ s。该阶段模型快速习得基本的冲突规避策略与速度调控策略。

\textbf{阶段II——渐进优化期（200--600轮）：}累计奖励波动上升至约$+1100\sim1300$，冲突次数缓慢下降至约$18\sim22$次/episode，奖励波动幅度$<5\%$。课程学习在该阶段触发多次难度进阶（$N$从11逐步增至19），每次进阶后观测到短暂（1--3轮）的性能回落，随后快速恢复并超越进阶前水平。

\textbf{阶段III——精细收敛期（600--1000轮）：}$N$固定为21，模型在高密度场景下进行策略微调。累计奖励稳定在$+1286.7 \pm 18.2$（最后50轮的均值），冲突次数收敛至$18.2 \pm 2.1$。

\begin{figure}[H]
    \centering
    \includegraphics[width=0.8\textwidth]{QQ图片20260504231524.jpg}
    \caption{训练收敛曲线——累计奖励（左上）、冲突次数（右上）、平均延迟（左下）与PCR/CFCR（右下）随训练轮次的变化。阴影区域表示5次独立训练的$\pm1$标准差。红色虚线标记课程学习难度进阶的时间点。}
    \label{fig:training}
\end{figure}

\subsection{综合性能对比}

$N=21$架eVTOL测试场景下的综合性能对比如表\ref{tab:comparison}所示（5次独立测试，均值$\pm$标准差）。

\begin{table}[H]
    \centering
    \caption{综合性能对比（21架eVTOL，5次独立测试）}
    \label{tab:comparison}
    \begin{tabular}{|c|c|c|c|c|}
        \hline
        \textbf{方法（类别）} & \textbf{CR$\uparrow$} & \textbf{TC$\downarrow$} & \textbf{AD(s)$\downarrow$} & \textbf{PCR(\%)$\uparrow$} \\
        \hline
        RR（规则） & $-1256.3\pm89.2$ & $42.6\pm5.3$ & $1.58\pm0.12$ & $89.2\pm3.5$ \\
        FCFS（规则） & $-1187.5\pm76.4$ & $38.9\pm4.7$ & $1.42\pm0.10$ & $91.5\pm2.8$ \\
        DQN（单RL） & $-892.5\pm65.3$ & $28.3\pm3.9$ & $1.21\pm0.09$ & $94.7\pm2.1$ \\
        独立A2C（独立MARL） & $-657.8\pm52.6$ & $19.7\pm2.8$ & $0.93\pm0.07$ & $96.3\pm1.7$ \\
        ORCA（工程安全） & $-198.7\pm35.8$ & $22.8\pm2.9$ & $0.76\pm0.06$ & $96.0\pm1.6$ \\
        MADDPG（CTDE） & $-215.6\pm41.2$ & $25.4\pm3.2$ & $0.87\pm0.06$ & $95.8\pm1.9$ \\
        COMA（CTDE） & $187.3\pm35.7$ & $22.1\pm2.9$ & $0.79\pm0.06$ & $97.1\pm1.5$ \\
        QMIX（CTDE） & $426.5\pm28.9$ & $23.1\pm2.7$ & $0.76\pm0.05$ & $97.5\pm1.3$ \\
        MAPPO（CTDE） & $892.4\pm23.5$ & $20.3\pm2.4$ & $0.65\pm0.04$ & $98.2\pm1.1$ \\
        \textbf{本文方法} & $\mathbf{1286.7\pm18.2}$ & $\mathbf{18.2\pm2.1}$ & $\mathbf{0.57\pm0.03}$ & $\mathbf{99.1\pm0.8}$ \\
        \hline
    \end{tabular}
\end{table}

从表\ref{tab:comparison}可得出以下层次化分析：

\textbf{规则方法与学习方法的系统性差距。}本文方法相比RR：TC降低57.2\%（42.6$\to$18.2），AD缩短63.9\%（1.58$\to$0.57 s），PCR提升9.9个百分点。规则方法虽有计算成本极低（$<1$ ms/步）的优势，但其性能上限受限于固定调度逻辑无法感知飞行器间相对运动状态。

\textbf{工程安全方法的性能定位。}ORCA的TC为22.8，显著优于规则方法（RR: 42.6，FCFS: 38.9），验证了分布式几何避障的有效性。但其CFCR仅约50\%（见表\ref{tab:safety}），低于本文方法（68.4\%），主要原因在于ORCA仅基于当前几何关系决策，缺乏对未来多步轨迹的预测能力，且在$N=21$的高密度场景中速度空间的可行解可能出现退化。

\textbf{与MAPPO的差异——数值温和但结构重要。}本文方法相比MAPPO的TC改进为10.3\%（20.3$\to$18.2），AD缩短12.3\%（0.65$\to$0.57 s）。改进幅度在数值上是温和的——MAPPO本身已具备相当的冲突规避能力。本文方法相较MAPPO的增量主要来自DyGAT的三个组件协同：空间注意力比全连接网络更精准地聚焦高风险邻域飞行器；时序注意力利用历史轨迹预判冲突趋势；物理量驱动的动态边权重为注意力提供了有价值的先验引导。

\subsubsection{安全指标专项对比}

\begin{table}[H]
    \centering
    \caption{安全指标专项对比（21架eVTOL，5次独立测试）}
    \label{tab:safety}
    \begin{tabular}{|c|c|c|c|c|}
        \hline
        \textbf{方法} & \textbf{CFCR(\%)$\uparrow$} & \textbf{MCAE$\downarrow$} & \textbf{MSVR(\%)$\downarrow$} & \textbf{PCR(\%)$\uparrow$} \\
        \hline
        RR & $12.3\pm4.1$ & $2.03\pm0.25$ & $8.7\pm1.2$ & $89.2\pm3.5$ \\
        FCFS & $18.7\pm3.6$ & $1.85\pm0.22$ & $7.4\pm1.0$ & $91.5\pm2.8$ \\
        ORCA & $50.2\pm3.0$ & $1.09\pm0.14$ & $3.8\pm0.6$ & $96.0\pm1.6$ \\
        DQN & $37.5\pm3.2$ & $1.35\pm0.19$ & $5.1\pm0.8$ & $94.7\pm2.1$ \\
        独立A2C & $48.2\pm2.8$ & $0.94\pm0.13$ & $3.6\pm0.6$ & $96.3\pm1.7$ \\
        MADDPG & $42.6\pm3.0$ & $1.21\pm0.15$ & $4.2\pm0.7$ & $95.8\pm1.9$ \\
        COMA & $51.3\pm2.7$ & $1.05\pm0.14$ & $3.8\pm0.6$ & $97.1\pm1.5$ \\
        QMIX & $53.7\pm2.5$ & $1.10\pm0.13$ & $3.5\pm0.5$ & $97.5\pm1.3$ \\
        MAPPO & $57.8\pm2.2$ & $0.97\pm0.11$ & $3.0\pm0.4$ & $98.2\pm1.1$ \\
        \textbf{本文} & $\mathbf{68.4\pm1.8}$ & $\mathbf{0.87\pm0.10}$ & $\mathbf{2.4\pm0.3}$ & $\mathbf{99.1\pm0.8}$ \\
        \hline
    \end{tabular}
\end{table}

表\ref{tab:safety}揭示了PCR与CFCR之间的巨大鸿沟——本文方法PCR $=99.1\%$但CFCR仅$68.4\%$，相差约30个百分点。MAPPO的PCR（98.2\%）与CFCR（57.8\%）之间也存在约40个百分点的差距，表明这一问题是CTDE-MARL方法的共性挑战。MCAE $=0.87$表明平均每架飞行器每episode经历不到1次冲突，较MAPPO的0.97次降低10.3\%。MSVR $=2.4\%$表明仅2.4\%的仿真时间步存在安全距离违反。

\subsection{冲突严重程度分级}

为分析CFCR $=68.4\%$中剩余约31\%存在冲突的episode的严重程度，将安全距离违反按$d_{\min}/d_{\text{safety}}$比值分为三级：轻微（$0.8 \leq d_{\min}/d_{\text{safety}} < 1.0$）、中等（$0.5 \leq d_{\min}/d_{\text{safety}} < 0.8$）、严重/NMAC（$d_{\min}/d_{\text{safety}} < 0.5$，对应$d_{\min} < 42.5$ m）。分级统计（50次测试episode）表明约80\%的违反事件属于轻微级别，NMAC级别事件占比低于1\%。

\subsection{统计显著性检验}

对本文方法与MAPPO在四项关键指标上的性能差异进行Welch双样本$t$检验（双侧，$\alpha=0.05$，$n_1=n_2=5$，不等方差），结果如表\ref{tab:ttest}所示。

\begin{table}[H]
    \centering
    \caption{本文 vs.\ MAPPO统计检验结果}
    \label{tab:ttest}
    \begin{tabular}{|c|c|c|c|c|c|}
        \hline
        \textbf{指标} & \textbf{均值差} & \textbf{$t$(df=8)} & \textbf{$p$} & \textbf{Cohen's $d$} & \textbf{结论} \\
        \hline
        TC & $-2.1$ & $1.47$ & $0.180$ & $0.93$ & 不显著 \\
        CFCR & $+10.6\%$ & $8.34$ & $<0.001$ & $5.27$ & 高度显著 \\
        MCAE & $-0.10$ & $1.50$ & $0.172$ & $0.95$ & 不显著 \\
        MSVR & $-0.6\%$ & $2.68$ & $0.029$ & $1.69$ & 显著 \\
        \hline
    \end{tabular}
\end{table}

TC的差异具有大效应量（Cohen's $d=0.93$），但由于$n=5$的小样本量限制，统计效力（power）仅约38\%——这意味着即使真实差异存在，当前样本量也只有约38\%的概率在$\alpha=0.05$水平上检测到。要达到80\%统计效力需每组$n \geq 20$。因此，TC的"不显著"恰当解读为"在$n=5$条件下未检测到统计显著差异"，而非"确认两组真实均值相等"。相比之下，CFCR的效应量极大（$d=5.27$），即使$n=5$也达到$>99\%$的统计效力，结论高度稳健。

本文方法在CFCR和MSVR上的统计显著改善——而非TC——构成了核心统计贡献。

\subsection{消融实验}

表\ref{tab:ablation}报告了各核心模块的独立性能贡献。

\begin{table}[H]
    \centering
    \caption{消融实验结果（21架eVTOL，5次独立测试，均值$\pm$标准差）}
    \label{tab:ablation}
    \resizebox{0.92\textwidth}{!}{%
    \begin{tabular}{|c|c|c|c|c|}
        \hline
        \textbf{模型配置} & \textbf{CR$\uparrow$} & \textbf{TC$\downarrow$} & \textbf{AD(s)$\downarrow$} & \textbf{PCR(\%)$\uparrow$} \\
        \hline
        完整模型 & $\mathbf{1286.7\pm18.2}$ & $\mathbf{18.2\pm2.1}$ & $\mathbf{0.57\pm0.03}$ & $\mathbf{99.1\pm0.8}$ \\
        \hline
        移除DyGAT（替换为标准GAT） & $215.3\pm32.6$ & $49.8\pm4.5$ & $0.79\pm0.05$ & $94.3\pm1.9$ \\
        \hline
        移除课程学习（直接$N=21$训练） & $587.2\pm27.4$ & $28.4\pm3.1$ & $0.68\pm0.04$ & $97.2\pm1.4$ \\
        \hline
        移除冲突损失加权（$\lambda \equiv 1.0$） & $763.5\pm24.1$ & $24.6\pm2.8$ & $0.65\pm0.04$ & $97.8\pm1.2$ \\
        \hline
        同时移除课程学习与冲突加权 & $128.6\pm38.5$ & $35.7\pm3.9$ & $0.82\pm0.06$ & $95.6\pm1.7$ \\
        \hline
    \end{tabular}%
    }
\end{table}

DyGAT是性能贡献最大的单一模块：移除后TC从18.2增至49.8（相对增幅173.6\%），CR从1286.7骤降至215.3。当前消融将DyGAT整体替换为标准GAT（仅保留空间注意力，移除时序注意力和动态边权重），尚未分解各子模块的独立贡献。

课程学习与冲突损失加权形成互补协同：单独移除分别导致TC增加56.0\%和35.2\%，同时移除增加96.2\%——两者效应之和与联合效应基本一致，表明两种机制的作用路径大部分相互独立。

\subsection{密度鲁棒性与敏感性分析}

\textbf{密度鲁棒性。}在$N \in \{5, 10, 15, 21\}$四种密度下（见图\ref{fig:density}），本文方法的TC增长斜率（$\approx0.75$次/架）显著低于MAPPO（$\approx1.15$次/架）和RR（$\approx2.4$次/架）。DyGAT在密度升高时优势愈发显著：飞行器越多，邻域内飞行器对越多，精准区分不同飞行器对的冲突风险等级越关键。

\begin{figure}[H]
    \centering
    \includegraphics[width=0.65\textwidth]{QQ图片20260524171218.png}
    \caption{不同飞行器密度（$N=5, 10, 15, 21$）下的冲突次数对比。误差条为$\pm1$标准差（5次独立测试）。}
    \label{fig:density}
\end{figure}

\textbf{超参数敏感性。}在$N=21$场景下测试了三类关键参数的影响（见图\ref{fig:sensitivity}）：$d_{\text{base}}$从40 m增至70 m时，TC下降34.7\%但AD上升21.2\%，体现经典的安全-效率权衡；风扰强度从0增至0.3 m/s时TC上升约28\%；传感器噪声$\sigma$从0增至1.5 m时TC上升约15\%——影响最小，LSTM编码器对历史序列的隐式平滑作用有效过滤了高频观测噪声。

\begin{figure}[H]
    \centering
    \includegraphics[width=0.8\textwidth]{QQ图片20260524162300.png}
    \caption{超参数敏感性分析——改变基础安全距离（左）、风扰强度（中）与传感器噪声（右）时冲突次数与平均延迟的变化。红色虚线标记默认参数值。误差条为$\pm1$标准差。}
    \label{fig:sensitivity}
\end{figure}

\subsection{计算成本分析}

本文方法在$N=21$时平均单步推理时间为$4.5$ ms（NVIDIA RTX 4060），远低于仿真步长500 ms，实时性裕度约$111\times$（见图\ref{fig:computation}）。计算时间从$N=5$时的$1.2$ ms增长至$N=21$时的$4.5$ ms，增长斜率约$0.21$ ms/架，在所有CTDE-MARL方法中最低。

\begin{figure}[H]
    \centering
    \includegraphics[width=0.8\textwidth]{QQ图片20260524163903.jpg}
    \caption{不同飞行器数量下的平均单步推理计算时间对比（NVIDIA RTX 4060 GPU）。}
    \label{fig:computation}
\end{figure}

%==========================================================================
% 5. 讨论
%==========================================================================
\section{讨论}

\subsection{与MS-GRL的比较}

Li等\cite{ref10}提出的MS-GRL与本文目标最为接近。两者在三个层面存在本质差异：

\textbf{（1）算法选择。}MS-GRL采用Double DQN处理离散动作空间（速度档位选择），训练稳定但控制粒度粗糙。本文采用PPO处理连续加速度控制，适合eVTOL的精细化速度调控需求——在连续动作空间中具有更好的稳定性。

\textbf{（2）图结构设计。}MS-GRL的多尺度GAT基于固定距离阈值（100 m、300 m、500 m），无法区分"距离相同但相对运动状态不同"的飞行器对的风险差异。以具体数值说明：两对距离均为200 m的飞行器，若一对同向等速飞行（$w_{ij} \approx 1.0$），另一对相向高速靠近（$w_{ij} \approx 0.015$），本文DyGAT对后者的注意力权重仅为前者的约1/67——固定阈值图完全无法捕捉这一差异。

\textbf{（3）场景规模与针对性。}MS-GRL面向100架UAV通用场景，本文聚焦21架eVTOL交叉路口的特定运营场景，在安全约束（动态安全距离）、物理建模（运动学约束）与运营逻辑上更具针对性。

两者在图增强MARL技术路线上形成互补。

\subsection{核心发现与机制解读}

\textbf{CFCR作为核心安全指标。}本文最重要的实证发现是PCR（99.1\%）与CFCR（68.4\%）之间约30个百分点的差距——这一差距在MAPPO方法中更大（PCR 98.2\% vs.\ CFCR 57.8\%，差距约40个百分点）。这表明仅报告PCR可能系统性高估CTDE-MARL方法的安全性能。CFCR的引入为低空交通管控的MARL评估提供了一个更诚实的基准。

\textbf{DyGAT的三种注意力机制协同。}消融实验中DyGAT移除导致TC增加173.6\%，远超其他任何模块。但当前消融未分解空间注意力、时序注意力与动态边权重的独立贡献——三者各自的增量效应和潜在的交互效应有待进一步研究。

\textbf{统计检验力的差异化图景。}TC的大效应量（$d=0.93$）与统计不显著（$p=0.180$）的组合，揭示了$n=5$小样本量下统计推断的根本限制——这种"大效应量但统计不显著"的模式在安全攸关的AI系统评估中应被更加谨慎地对待。

\subsection{局限性}

本文方法存在以下主要局限性：

\textbf{（1）二维场景限制。}当前仿真为二维单交叉路口场景。真实低空交通涉及三维高度层（含垂直起降与巡航高度切换）、多路口网络拓扑、起降场容量约束等维度。

\textbf{（2）统计样本量。}$n=5$的独立测试次数对中等效应量指标（如TC的$d=0.93$）的统计效力不足（约38\%）。扩大样本量至$n \geq 20$以提供更可靠的统计推断是必要的后续工作。

\textbf{（3）风扰与噪声模型简化。}当前采用均匀随机风扰和高斯白噪声，均为时间独立的简化模型。真实大气边界层风场具有空间相关性和时间连续性（Dryden/Kaimal谱），城市环境还存在建筑尾流等复杂效应。

\textbf{（4）通信假设。}CTDE训练假设全局状态可无延迟获取。实际UTM部署中通信链路存在延迟（50--200 ms）和丢包约束。

\textbf{（5）基线覆盖。}MPC和CBF两种具有形式化安全保证的方法尚未纳入实验对比。ORCA的关键参数已进行了针对本文场景的调优，但未对所有基线方法进行系统性的超参数搜索。

\textbf{（6）eVTOL动力学参数。}加速度取值参考了学术文献中的典型建模假设\cite{ref25}，当前主流制造商并未公开发布精确的最大水平加速度规格。在更精确的制造商数据可用时应进行参数敏感性验证。

%==========================================================================
% 6. 结论
%==========================================================================
\section{结论}

本文针对城市低空十字交叉路口的eVTOL协同管控问题，提出了一种融合DyGAT的耦合MARL模型。在21架eVTOL的二维交叉路口仿真场景中，本文方法在无冲突完成率（CFCR：68.4\% vs.\ MAPPO的57.8\%，$p<0.001$）和最小间隔违反率（MSVR：2.4\% vs.\ 3.0\%，$p=0.029$）上取得了统计显著的改善。总冲突次数的差异（TC：18.2 vs.\ 20.3，降幅10.3\%）具有大效应量（$d=0.93$），但在$n=5$的样本量下未达统计显著（$p=0.180$）——更大规模的重复实验是验证该效应的必要步骤。

本文的方法论贡献在于：（1）物理量驱动的动态边权重为图结构MARL中的时空耦合建模提供了一种有效方法；（2）区分表层与深层指标的评价框架揭示了PCR与CFCR之间的系统性差距；（3）对统计效力的诚实分析展示了当前实证证据的强度与局限。

后续工作方向包括：将场景扩展至三维多交叉路口网络；在更大计算资源下进行$n \geq 20$的独立重复实验；集成标准风谱模型（Dryden/Kaimal）替代当前简化风扰；引入通信延迟与丢包约束；与MPC、CBF等安全关键方法进行系统基准对比。

%==========================================================================
% 参考文献
%==========================================================================
\begin{thebibliography}{99}

\bibitem{ref1} Federal Aviation Administration. Urban air mobility (UAM): concept of operations v2.0[R]. Washington, DC: FAA, 2023.

\bibitem{ref2} MENON P K, SWERIDUK G D, BILIMORIA K D. New approach to modeling, analysis, and control of air traffic flow[J]. Journal of Guidance, Control, and Dynamics, 2004, 27(5): 737--744.

\bibitem{ref3} RICHARDS A, HOW J P. Aircraft trajectory planning with collision avoidance using mixed integer linear programming[C]//Proceedings of the 2002 American Control Conference. Piscataway: IEEE, 2002: 1936--1941.

\bibitem{ref4} VAN DEN BERG J, GUY S J, LIN M, et al. Reciprocal n-body collision avoidance[C]//Robotics Research: The 14th International Symposium ISRR (Springer Tracts in Advanced Robotics, Vol.\ 70). Berlin: Springer, 2011: 3--19.

\bibitem{ref5} AMES A D, COOGAN S, EGERSTEDT M, et al. Control barrier functions: theory and applications[C]//2019 18th European Control Conference (ECC). Piscataway: IEEE, 2019: 3420--3431.

\bibitem{ref6} AGOGINO A, TUMER K. Regulating air traffic flow with coupled agents[C]//Proceedings of the 7th International Conference on Autonomous Agents and Multiagent Systems (AAMAS 2008). Richland: IFAAMAS, 2008: 535--542.

\bibitem{ref7} AGOGINO A K, TUMER K. A multiagent approach to managing air traffic flow[J]. Autonomous Agents and Multi-Agent Systems, 2012, 24(1): 1--25.

\bibitem{ref8} BRITTAIN M, WEI P. Autonomous separation assurance in a high-density en route sector: a deep multi-agent reinforcement learning approach[C]//2019 IEEE Intelligent Transportation Systems Conference (ITSC). Piscataway: IEEE, 2019: 3256--3262.

\bibitem{ref9} GHOSH S, LAGUNA S, LIM S H, et al. A deep ensemble method for multi-agent reinforcement learning: a case study on air traffic control[C]//Proceedings of the 31st International Conference on Automated Planning and Scheduling (ICAPS 2021). Palo Alto: AAAI Press, 2021: 468--476.

\bibitem{ref10} LI Y, LI J, WANG J, ZHANG X, DING H, DU W. Multiscale-graph-enhanced reinforcement learning for conflict resolution in dense UAV networks[J]. IEEE Internet of Things Journal, 2025, 12(21): 44290--44303.

\bibitem{ref11} SPATHARIS C, BASTAS A, KRAVARIS T, et al. Hierarchical multiagent reinforcement learning schemes for air traffic management[J]. Neural Computing and Applications, 2021, 33(14): 8649--8671.

\bibitem{ref12} ASLANI M, MESGARI M S, WIERING M. Adaptive traffic signal control with actor-critic methods in a real-world traffic network with different traffic disruption events[J]. Transportation Research Part C: Emerging Technologies, 2017, 85: 732--752.

\bibitem{ref13} 翟子洋, 郝茹茹, 董世浩. 大规模智慧交通信号控制中的强化学习和深度强化学习方法综述[J]. 计算机应用研究, 2024, 41(6): 1618--1627.

\bibitem{ref14} TUMER K, AGOGINO A. Distributed agent-based air traffic flow management[C]//Proceedings of the 6th International Joint Conference on Autonomous Agents and Multiagent Systems (AAMAS 2007). Honolulu: ACM, 2007: 330--337.

\bibitem{ref15} 王迎. 城市交通信号控制深度强化学习算法研究[D]. 济南: 山东交通学院, 2024.

\bibitem{ref16} 赵晓华, 石建军, 李振龙, 赵国勇. 基于Q-learning和BP神经元网络的交叉口信号灯控制[J]. 公路交通科技, 2007, 24(7): 99--102.

\bibitem{ref17} 赖建辉. 基于深度强化学习的交通控制优化方法研究及实现[D]. 杭州: 浙江工业大学, 2020.

\bibitem{ref18} 张新武. 基于深度强化学习算法的交通信号控制优化研究[D]. 太原: 太原科技大学, 2024.

\bibitem{ref19} 承向军, 常歆识, 杨肇夏. 基于Q-学习的交通信号控制方法[J]. 系统工程理论与实践, 2006(8): 136--140.

\bibitem{ref20} 卢守峰, 张术, 刘喜敏. 平均排队长度差最小的单交叉口在线Q学习模型[J]. 公路交通科技, 2014, 31(11): 116--122.

\bibitem{ref21} 江波, 李彦冬, 刘威陇, 等. 应用深度Q学习的机场道口协同智能调速方法[J]. 航空计算技术, 2022, 52(1): 26--30.

\bibitem{ref22} 张春阳, 席雅琪. 人工智能在城市交通控制中的应用综述[J]. 信息系统工程, 2024(07): 33--36.

\bibitem{ref23} GHAVAMZADEH M, MAHADEVAN S, MAKAR R. Hierarchical multi-agent reinforcement learning[J]. Autonomous Agents and Multi-Agent Systems, 2006, 13(2): 197--229.

\bibitem{ref24} TUMER K, AGOGINO A. Improving air traffic management with a learning multiagent system[J]. IEEE Intelligent Systems, 2009, 24(1): 18--21.

\bibitem{ref25} BACCHINI A, CESTINO E. Electric VTOL configurations comparison[J]. Aerospace, 2019, 6(3): 26.

\bibitem{ref26} U.S. Department of Defense. Flying qualities of piloted airplanes: MIL-F-8785C[S]. Washington, DC: U.S. Department of Defense, 1980.

\bibitem{ref27} LIANG X, MA Y, FENG Y, et al. PTR-PPO: proximal policy optimization with prioritized trajectory replay[EB/OL]. arXiv:2112.03798, 2021.

\bibitem{ref28} CONG W, ZHANG S, CHEN J, et al. Do we really need complicated model architectures for temporal networks?[C]//Proceedings of the 11th International Conference on Learning Representations (ICLR 2023). Kigali: ICLR, 2023.

\bibitem{ref29} XIE Y B, GARDI A, SABATINI R. Reinforcement learning-based flow management techniques for urban air mobility and dense low-altitude air traffic operations[C]//2021 IEEE/AIAA 40th Digital Avionics Systems Conference (DASC). Piscataway: IEEE, 2021: 1--10.

\end{thebibliography}

\end{document}
